\documentclass{ctexart}
\usepackage[a4paper,top=2cm,bottom=2cm,left=3cm,right=3cm,marginparwidth=1.75cm]{geometry}

\usepackage{amsmath}
\usepackage{graphicx}
\usepackage[colorlinks=true, allcolors=blue]{hyperref}
\usepackage{CJKutf8}
\usepackage{CJKpunct}

\title{一元三次方程}
\author{笙}

\begin{document}
\tableofcontents
\maketitle


\begin{abstract}
补充内容.一元三次方程.
\end{abstract}

\section{概念}

\subsection{定义}

形如$y=ax^3+bx^2+cx+d (a\neq 0)$的函数。

在高中范围内，我们习惯性关注一元三次方程$ax^3+bx^2+cx+d=0(a\neq 0)$在实数域上的解。

\section{一元三次方程}

\subsection{韦达公式}

\begin{equation}
    x_1x_2x_3=-\frac{d}{a}
\end{equation}

\begin{equation}
    x_1x_2+x_1x_3+x_2x_3=\frac{c}{a}
\end{equation}

\begin{equation}
    x_1+x_2+x_3=-\frac{b}{a}
\end{equation}


\subsection{盛金公式}

对方程\textbf{$ax^3+bx^2+cx+d=0(a\neq 0)$}。

\bigskip

有\textbf{重根判别式}：

\begin{align}
    \begin{cases}
        A=b^2-3ac\\
        B=bc-9ad\\
        C=c^2-3bd
    \end{cases}    
\end{align}

\textbf{总判别式}

\begin{equation}
    \Delta=B^2-4AC
\end{equation}

\begin{table}[]
    \centering
    \begin{tabular}{|c|c|}
        \hline
        $\Delta \textgreater 0$ & 方程有一个实根和两个共轭副根 \\
        \hline
        $\Delta = 0$ & 方程有三个实根(含有一个二重根)\\
        \hline
        $\Delta \textless 0$ & 方程有三个互不相等的实根 \\ 
        \hline
    \end{tabular}
\end{table}

\subsection{$\Delta \textgreater 0$}

令$Y_1$、$Y_2$为：

\begin{equation}
    Y_{1,2}=Ab+3a(\frac{-B\pm \sqrt{B^2-4AC}}{2})
\end{equation}

实根：

\begin{equation}
    X_1=\frac{-b-(\sqrt[3]{Y_1}+\sqrt[3]{Y_2})}{3a}
\end{equation}

共轭复根：
\begin{equation}
    X_{2,3}=\frac{-b+\frac{1}{2}(\sqrt[3]{Y_1}+\sqrt[3]{Y_2})\pm \frac{\sqrt{3}}{2}(\sqrt[3]{Y_1}-\sqrt[3]{Y_2})i}{3a}
\end{equation}

\subsection{$\Delta = 0$}

令$K=\frac{B}{A}$。

单根：

\begin{equation}
    X_1=-\frac{b}{a}+K
\end{equation}

二重根：

\begin{equation}
    X_{2,3}=-\frac{K}{2}
\end{equation}

\subsection{$\Delta \textless 0$}

令$\theta$、$T$分别为：

\begin{equation}
    \theta=\arccos T
\end{equation}

\begin{equation}
    T=\frac{2Ab-3aB}{2\sqrt{A^3}}
\end{equation}

则根为：

\begin{equation}
    X_1=\frac{-b-2\sqrt{A}\cos\frac{\theta}{3}}{3a}
\end{equation}

\begin{equation}
    X_{2,3}=\frac{-b+\sqrt{A}(\cos\frac{\theta}{3}\pm \sqrt{3}\sin\frac{\theta}{3})}{3a}
\end{equation}

\end{document}